\documentclass[a4paper,10.5pt]{article}
\usepackage[T1]{fontenc}\usepackage[utf8]{inputenc}
\usepackage[italian]{babel}
\usepackage[top=2.2cm,bottom=2.2cm,left=2.35cm,right=2.35cm]{geometry}
\usepackage{hyperref}\usepackage{xcolor}\definecolor{accent}{HTML}{1F4E78}\hypersetup{colorlinks=true,urlcolor=accent}
\pagestyle{empty}\setlength{\parindent}{0pt}\setlength{\parskip}{7pt}
\begin{document}
\textbf{Prabhavit Gupta}\\
Venezia, Italia \textbar{} +39 351 610 7807 \textbar{} prabhavitg@gmail.com\\[10pt]
Gentile team di Mediatica Digital,

vi scrivo per candidarmi al ruolo di Mobile Test Automation Engineer. Negli ultimi tre anni ho lavorato su automazione dei test, scraping di dati e pipeline di qualità in ambienti Agile, ed è quel lato del ruolo che posso portare fin da subito, insieme a un interesse concreto ad approfondire l'automazione mobile con Appium.

Per Clubee, in Lussemburgo, ho fatto parte del team che integrava le API Clubee su tutto il frontend e mi sono occupato di unit testing e quality assurance delle nuove funzionalità, lavorando a stretto contatto con Git e GitHub per il versionamento del codice, in un contesto Agile con rilasci frequenti.

Per FlowSpace, un marketplace di spazi di coworking con sede a Londra, uso Playwright per automatizzare test end-to-end e per lo scraping e la verifica dei dati che alimentano la piattaforma, mentre gestisco l'intero flusso di lavoro su GitHub e Jira insieme al team di sviluppo.

Non ho ancora esperienza diretta con Appium, ma la mia base in Playwright e nella progettazione di suite di test automatizzati mi permette di imparare in fretta un nuovo framework di test mobile, ed è un'area che mi interessa sviluppare con continuità.

Ho status di soggiornante di lungo periodo UE in Italia e sono autorizzato a lavorare senza necessità di sponsorizzazione. Sono disponibile per lavoro full remote, con occasionali trasferte a Roma, e posso iniziare da subito.

Cordiali saluti,\\
Prabhavit Gupta
\end{document}
